% !TEX root = ../main.tex
\chapter{从研究问题到论文结构}
\label{chap:guide1}

\section{先写清楚问题，再安排章节}

开始写作时，可以先用一小段话说明研究对象、希望回答的问题、已有材料和当前限制。问题应当具体到能够判断哪些材料与它有关。例如，“研究写作工具”只是主题；“说明一份材料如何组织为可编辑章节，并记录仍需确认的信息”才给出了可执行的写作目标。这里的区别不是句子长短，而是读者能否据此理解讨论范围。尚未取得的材料应单独列出，不要把预计完成的工作写成已经得到的结果。

问题陈述还需要与研究结论相互对应。如果开头讨论的是引用信息如何追溯，结尾就应回答来源是否明确、哪些信息仍不完整，而不能突然转向工具是否提升科研创新。写作过程中可以修改问题，但每次调整都要检查摘要、章节标题和结论是否同步。为同一个概念保持相同名称，通常比增加更多术语更有帮助。本项目配套的章节组织说明给出了问题、材料和章节之间的记录方法。\cite{guide1}

\section{为材料建立可追溯的清单}

在正文之外维护一份材料清单，记录每项材料的名称、来源、版本、准备使用的位置和待核对事项。自己的记录也应注明产生时间与用途；引用他人材料时，应保留能够定位来源的信息。材料清单的目的不是积累文件数量，而是让每个论点都能回到具体证据。如果两个文件来自不同修改阶段，先判断哪一个是本次写作的依据，再合并内容，避免把相互矛盾的数值拼在同一段中。

可以将“材料已经提供”“材料已经核对”和“内容已经写入”作为不同状态。收到一张表格，并不等于表中指标的定义已经清楚；把文字放进章节，也不等于完成了对照审阅。当材料缺少单位、条件或来源时，应记录缺项与影响范围。这里不要求使用专门的软件，一个包含固定字段的普通文本表格就可以承担这些工作。涉及未公开数据时，公开示例应另用独立的教学材料。

\section{用章节职责约束写作范围}

章节提纲需要说明每一章解决什么问题，而不只是列出几个宽泛标题。表~\ref{tab:guide1} 给出一种通用的职责划分；实际论文应根据研究类型调整，不必机械套用。理论研究、实验研究和工程实践的材料不同，章节之间的关系也可能不同。判断提纲是否清楚，可以尝试用一句话描述相邻两章的衔接：前一章给出了什么，后一章为什么需要接着讨论。

\begin{table}[htbp]
  \centering
  \small
  \renewcommand{\arraystretch}{1.25}
  \caption{常见章节职责与检查问题}
  \label{tab:guide1}
  \begin{tabular}{@{}>{\raggedright\arraybackslash}p{.2\linewidth}>{\raggedright\arraybackslash}p{.52\linewidth}>{\raggedright\arraybackslash}p{.2\linewidth}@{}}
    \toprule
    章节角色 & 主要职责 & 自查问题 \\
    \midrule
    绪论 & 明确问题、背景与范围 & 读者能否说出本文要回答什么？ \\
    方法 & 交代步骤、变量与适用条件 & 信息是否足以理解操作过程？ \\
    结果与讨论 & 展示观察并解释其边界 & 结论能否回到具体材料？ \\
    总结 & 回答问题并列出局限 & 是否引入了正文未说明的新结论？ \\
    \bottomrule
  \end{tabular}
\end{table}

摘要应在正文结构稳定后回头校对，至少交代讨论对象、采用的方式、主要内容和适用边界。中英文摘要可以采用各自自然的表达，但研究对象、数量、关键条件和结论强度应一致。章节规划并不意味着从第一章顺序写到最后一章；作者可以先整理证据最充分的部分，再补上连接这些部分的解释。修改顺序可以灵活，论证关系必须清楚。
